\documentclass{article}
\usepackage[T1]{fontenc}
\usepackage{lmodern}
\usepackage{graphicx}
\usepackage{amsmath,amssymb}
\usepackage[a4paper,margin=2cm]{geometry}
\usepackage[absolute,overlay]{textpos}
\setlength{\TPHorizModule}{1mm}
\setlength{\TPVertModule}{1mm}
\usepackage[indonesian]{babel}
\usepackage{hyperref}
\setlength{\emergencystretch}{2em}
% unit: Halaman 2

\begin{document}

% === Halaman 2 (python/native) ===

Referensi:

1. Andreas Steffen, Elliptic Curve Cryptography, Zürcher Hochschule 

Winterthur.

2. Debdeep Mukhopadhyay, Elliptic Curve Cryptography , Dept of 

Computer Sc and Engg IIT Madras.

3. Anoop MS , Elliptic Curve Cryptography,  an Implementation Guide

2

Rinaldi Munir/II4021 Kriptografi

\end{document}
